\section{Method}
\label{sec:method}

We investigate where action-relevant representations originate within a vision-language-action (VLA) model by combining representational similarity analysis, functional probing, and kernel alignment across a factorial design of training and instruction conditions.

\subsection{Model and Data}
\label{sec:model_data}

\paragraph{Architecture.}
We study OpenVLA-7B~\cite{kim2024_openvla,karamcheti2024}, a VLA model comprising four components: (1)~a SigLIP-SO400M vision encoder~\cite{zhai2023_siglip} that maps RGB observations to patch-level visual tokens, (2)~an MLP projector that maps vision tokens into the language model's embedding space, (3)~a Llama-2 7B language backbone (32 transformer layers) that processes the concatenated sequence of projected image tokens and tokenized language instructions, and (4)~an action head that decodes 7-dimensional continuous actions (3D translation, 3D rotation, 1D gripper) from the final hidden states.

\paragraph{Training data.}
We fine-tune OpenVLA on the LIBERO-10 split~\cite{liu2023_libero}, which contains 10 tabletop manipulation tasks with 50 human demonstrations each (500 demonstrations total).
We evaluate generalization on LIBERO-Goal, an out-of-distribution benchmark where episodes share initial object configurations but differ in goal specifications.

\paragraph{Untrained control.}
To isolate the LLM's learned contribution from the vision encoder's pretrained geometry, we construct an untrained control by calling \texttt{from\_config()} to reinitialize all LLM parameters (6.74B weights drawn from random initialization) while preserving the frozen SigLIP encoder (0.80B pretrained weights).
The untrained model processes identical visual inputs through an identical vision encoder but lacks any learned language or action structure in its LLM layers.

\subsection{Representation Extraction}
\label{sec:rep_extraction}

We extract representations from five sources spanning the VLA processing pipeline and two external baselines:

\begin{enumerate}[label=(\alph*),nosep,leftmargin=*]
    \item \textbf{SigLIP+Projector} --- the vision encoder output after MLP projection into the LLM's input space, before any LLM processing. This captures the representation that the LLM receives as input.
    \item \textbf{LLM image\_mean} --- the mean activation across all image token positions at each of the 32 LLM layers. This captures how the LLM transforms distributed visual information through its depth.
    \item \textbf{LLM last\_preaction} --- the activation at the final token position before action decoding (the readout token) at each LLM layer. This captures the representation from which actions are directly decoded.
    \item \textbf{DINOv2 ViT-L/14}~\cite{oquab2024_dinov2} \texttt{cls\_token} --- a vision-only baseline pretrained with self-supervised objectives and no language supervision.
    \item \textbf{CLIP ViT-L/14}~\cite{radford2021_clip} \texttt{patch\_mean} --- a vision-language baseline pretrained on image-text pairs, serving as a generalization control for VL pretraining effects independent of the VLA pipeline.
\end{enumerate}

Sources (a)--(c) probe representations within the VLA model itself, while (d) and (e) provide external references with no VLA training signal.

\paragraph{Factorial design.}
We cross two factors in a $2 \times 2$ design: \emph{training status} (trained vs.\ untrained LLM) and \emph{instruction presence} (with vs.\ without language instruction).
This yields four conditions per LLM representation source, enabling us to disentangle the contributions of LLM training and language grounding to action-relevant structure.
External baselines (d) and (e) are evaluated once, as they are independent of both factors.

\subsection{Representational Similarity Analysis}
\label{sec:rsa}

We adopt Representational Similarity Analysis (RSA)~\cite{kriegeskorte2008} to quantify the geometric alignment between neural representations and the action space.

\paragraph{Representational dissimilarity matrices.}
For a set of $n$ observations, we compute an $n \times n$ representational dissimilarity matrix (RDM) for each representation source using Euclidean distance between representation vectors.
We similarly construct an action RDM from the 7-dimensional action vectors.

\paragraph{RSA correlation.}
We measure alignment between a representation RDM and the action RDM via Spearman rank correlation $\rho$ between their upper-triangular elements.
Higher $\rho$ indicates that observations with similar representations also have similar actions, i.e., the representation geometry mirrors the action space geometry.

\paragraph{Statistical inference.}
We compute 95\% confidence intervals from 10{,}000 bootstrap resamples of observation pairs for all pairwise comparisons.
RSA is computed per task and aggregated across the 10 tasks.

\subsection{Functional Action Probes}
\label{sec:probes}

RSA measures geometric alignment but does not assess whether action information is linearly decodable from representations.
We complement RSA with functional probing~\cite{alain2017} to measure the action information accessible to a lightweight decoder.

\paragraph{Probe architecture.}
We train two-layer MLP probes (Linear $\to$ ReLU $\to$ Linear) mapping from the representation dimensionality to a 256-dimensional hidden layer and then to the 7-dimensional action space.

\paragraph{Training protocol.}
Probes are optimized with Adam (learning rate $10^{-3}$, batch size 256) for 50 epochs with early stopping (patience 5 epochs on validation loss).
Representations are frozen throughout; only probe parameters are updated.
We use per-task 80/20 train/validation splits and report mean $R^2$ across the 7 action dimensions as the primary metric.

\paragraph{Action-delta probes.}
In addition to predicting absolute actions $\mathbf{a}_t$, we train probes to predict action deltas $\Delta\mathbf{a} = \mathbf{a}_{t+1} - \mathbf{a}_t$, which capture the dynamic component of action structure independent of the current action magnitude.

\paragraph{Statistical comparison.}
We use bootstrap paired comparison (10{,}000 resamples) to test whether $R^2$ differences between representation sources are statistically significant.

\subsection{Centered Kernel Alignment}
\label{sec:cka}

RSA captures local geometric alignment---whether pairwise distances in representation space mirror those in action space.
To assess global structural preservation, we additionally employ Centered Kernel Alignment (CKA)~\cite{kornblith2019} in two variants.

\paragraph{Linear CKA.}
Linear CKA computes the normalized Hilbert-Schmidt Independence Criterion between linear kernel matrices of two representation sets.
It is invariant to orthogonal transformations and isotropic scaling, making it sensitive to global structure rather than local distance preservation.

\paragraph{RDM kernel CKA.}
We also compute CKA using the RDMs themselves as kernel matrices, providing a bridge between the local sensitivity of RSA and the global perspective of linear CKA.

\paragraph{Dimensionality caveat.}
When comparing high-dimensional representations (e.g., 4096-D LLM activations) with low-dimensional action vectors (7-D), linear CKA may be dominated by the few directions that span the action space, potentially inflating similarity estimates.
We report both CKA variants alongside RSA to ensure that conclusions are robust to metric choice.
